\input{../tex-inputs/latex-leseansicht-vorspann}

\section[Olga und Arthur Schnitzler an Gustav Schwarzkopf, 21. 5. 1910]{L04372 Olga und Arthur Schnitzler an Gustav Schwarzkopf, 21. 5. 1910}
\nopagebreak\mylabel{L04372v}
\rehead{ }\normalsize\beginnumbering\briefempfaengerindex{Schwarzkopf, Gustav@\textsc{Schwarzkopf, Gustav}!zzzSchnitzler, Arthur@\emph{von Arthur Schnitzler}!1910-05-212@{21. 5. 1910}|(be}\briefempfaengerindex{Schwarzkopf, Gustav@\textsc{Schwarzkopf, Gustav}!zzzSchnitzler, Olga@\emph{von Olga Schnitzler}!1910-05-212@{21. 5. 1910}|(be}
\toendnotes[C]{\smallbreak\pagebreak[2]}
\correspDesc{Versand  durch Olga Schnitzler, Arthur Schnitzler am 21. 5. 1910 in Morschach
\newline{}Übermittlung  am 21. 5. 1910 in Morschach
\newline{}Erhalt  durch Gustav Schwarzkopf im Zeitraum [22. 5. 1910 – 26. 5. 1910?] in Tiefer Graben 23}\toendnotes[C]{\smallbreak}
\Standort{DLA, A:Schnitzler, HS.1985.1.1897.}
\physDesc{Bildpostkarte, 258 Zeichen
\newline{}Handschrift Olga Schnitzler: Bleistift, lateinische Kurrent
\newline{}Handschrift Arthur Schnitzler: Bleistift, lateinische Kurrent
\newline{}Versand: Stempel: »\nobreak{}\oindex{Morschach@\textbf{Morschach}|pwk}Morschach, 21. V. 10\nobreak{}«.  }\toendnotes[C]{\smallbreak}\pstart{}{\pb}Herrn\pend{}\pstart{}Gustav Schwarzkopf\pend{}\pstart{}Wien I\oindex{I., Innere Stadt@\textbf{I., Innere Stadt}|pw}\pend{}\pstart{}Tiefer Graben 23\oindex{Wien@\textbf{Wien}!I., Innere Stadt@\textbf{I., Innere Stadt}!Tiefer Graben 23@\textbf{Tiefer Graben 23}|pw}.\pend{}{\bigskip}
\pstart
           \centering{}{\pb}\textcolor{gray}{\textbf{Axenstein\oindex{Axenstein@\textbf{Axenstein}|pw}, vom See\oindex{Vierwaldstättersee@\textbf{Vierwaldstättersee}|pwv} aus gesehen}}\pend
           \vspace{1em}
\pstart
           {\pb}Sonntag 21. Mai 1910.\pend
           \vspace{0.5em}
\pstart
           {\pb}Herzliche Grüsse!
      Es ist hier so schön, und 
      heiss, und die Berge sind
      voll Schnee, und man
      kann’s nicht beschreiben.
      (Wenigstens \uline{ich} nicht!)\pend
           \pstart \spacefill\mbox{Olga}.\pend{}\selectlanguage{ngerman}\vspace{1em}
\pstart
           \noindent{}{[}hs. A. Schnitzler:{]} Glücklicherweise 
      muß man nicht. Herzlichst\pend
           \pstart Ihr
               \spacefill\mbox{A.}\pend{}\selectlanguage{ngerman}\endnumbering\briefempfaengerindex{Schwarzkopf, Gustav@\textsc{Schwarzkopf, Gustav}!zzzSchnitzler, Arthur@\emph{von Arthur Schnitzler}!1910-05-212@{21. 5. 1910}|)be}\briefempfaengerindex{Schwarzkopf, Gustav@\textsc{Schwarzkopf, Gustav}!zzzSchnitzler, Olga@\emph{von Olga Schnitzler}!1910-05-212@{21. 5. 1910}|)be}\mylabel{L04372h}
\begin{anhang}
\end{anhang}\newcommand{\dateiname}{L04372}\newcommand{\titel}{Olga und Arthur Schnitzler an Gustav Schwarzkopf, 21. 5. 1910}\newcommand{\editorInnen}{Herausgegeben von Jahnke, SelmaMüller, Martin Anton}\input{../tex-inputs/latex-leseansicht-abspann}
